% GENERATED FILE. Edit canonical lessons, then run npm run generate:books.
% canonical-source-hash: fnv1a64:7b8d241542d3d162
% canonical-lessons: MW-C200-taraju, MW-C200-ser, MW-C200-kilo, MW-C200-pav, MW-C200-litar

\chapter{Scales and Weights}
\label{ch:mw-scales-and-weights}
\hlchaptermodality{200}

\begin{chapteropening}
\textbf{By the end of this chapter:} \emph{I can ask for a weight of something.}

Complete the last lesson of chapter 200: I can ask for a weight of something.
\end{chapteropening}

\section[\texorpdfstring{\mw{तराजू} (tarājū)}{तराजू (tarājū)}]{\mw{तराजू} (tarājū) — a pair of scales}

\label{lesson:MW-C200-taraju}



\begin{quote}
Before the new one: say the Marwadi for a reckoning, then the Marwadi for to write down.
\end{quote}

\subsection*{You'll want to know: \mw{तराजू}}
\textbf{\mw{तराजू}} — \emph{tarājū} — \textquotedblleft{}a pair of scales\textquotedblright{}.

\begin{quote}
\mw{तराजू} \mw{दुकान} \mw{मांय} \mw{है।}
\end{quote}

\begin{quote}
\emph{tarājū dukān mā̃y hai.}
\end{quote}

\begin{quote}
The scales are in the shop.
\end{quote}

A masculine noun.

\begin{cousinweb}
\textbf{\mw{तराजू}} (\emph{tarājū}) is usually traced to Persian \emph{tarāzū}, \textquotedblleft{}balance, scales\textquotedblright{}.
\end{cousinweb}

\begin{grammarlens}[title={What you've built}]
One more word for reading notices and messages.
\end{grammarlens}

\subsection*{Your turn}
\begin{itemize}
\raggedright
  \item \emph{Say it:} \emph{tarājū}
  \item \emph{Say it:} \emph{tarājū}, once more
  \item \emph{Say it:} \emph{tarājū dukān mā̃y hai.} — The scales are in the shop.
  \item \emph{Recall:} say the Marwadi for a reckoning, then the Marwadi for to write down, then say \emph{tarājū} again
\end{itemize}

\begin{tcolorbox}[breakable,colback=teal!4,colframe=teal!35!black,title={Before you move on}]
What does \mw{तराजू} mean? (\textquotedblleft{}A pair of scales\textquotedblright{}.) Say it once more.
\end{tcolorbox}
\section[\texorpdfstring{\mw{सेर} (ser)}{सेर (ser)}]{\mw{सेर} (ser) — a seer (a measure of weight)}

\label{lesson:MW-C200-ser}



\begin{quote}
Before the new one: say the Marwadi for to write down, then the Marwadi for scales.
\end{quote}

\subsection*{You'll want to know: \mw{सेर}}
\textbf{\mw{सेर}} — \emph{ser} — \textquotedblleft{}a seer (a measure of weight)\textquotedblright{}.

\begin{quote}
\mw{एक} \mw{सेर} \mw{आटो} \mw{दो।}
\end{quote}

\begin{quote}
\emph{ek ser āṭo do.}
\end{quote}

\begin{quote}
Give me a seer of flour.
\end{quote}

A masculine noun.

\begin{cousinweb}
\textbf{\mw{सेर}} (\emph{ser}) is inherited from a reconstructed \emph{*satēra}, \textquotedblleft{}a measure of weight\textquotedblright{}; CDIAL lists Old Marwari \emph{sera} (CDIAL 13106).
\end{cousinweb}

\begin{grammarlens}[title={What you've built}]
One more word for reading notices and messages.
\end{grammarlens}

\subsection*{Your turn}
\begin{itemize}
\raggedright
  \item \emph{Say it:} \emph{ser}
  \item \emph{Say it:} \emph{ser}, once more
  \item \emph{Say it:} \emph{ek ser āṭo do.} — Give me a seer of flour.
  \item \emph{Recall:} say the Marwadi for to write down, then the Marwadi for scales, then say \emph{ser} again
\end{itemize}

\begin{tcolorbox}[breakable,colback=teal!4,colframe=teal!35!black,title={Before you move on}]
What does \mw{सेर} mean? (\textquotedblleft{}A seer (a measure of weight)\textquotedblright{}.) Say it once more.
\end{tcolorbox}
\section[\texorpdfstring{\mw{किलो} (kilo)}{किलो (kilo)}]{\mw{किलो} (kilo) — a kilo}

\label{lesson:MW-C200-kilo}



\begin{quote}
Before the new one: say the Marwadi for scales, then the Marwadi for a seer.
\end{quote}

\subsection*{You'll want to know: \mw{किलो}}
\textbf{\mw{किलो}} — \emph{kilo} — \textquotedblleft{}a kilo\textquotedblright{}.

\begin{quote}
\mw{दो} \mw{किलो} \mw{चावळ} \mw{लावो।}
\end{quote}

\begin{quote}
\emph{do kilo chāvaḷ lāvo.}
\end{quote}

\begin{quote}
Bring two kilos of rice.
\end{quote}

A masculine noun.

\begin{cousinweb}
\textbf{\mw{किलो}} (\emph{kilo}) is borrowed from English \emph{kilo}, short for kilogram.
\end{cousinweb}

\begin{grammarlens}[title={What you've built}]
One more word for reading notices and messages.
\end{grammarlens}

\subsection*{Your turn}
\begin{itemize}
\raggedright
  \item \emph{Say it:} \emph{kilo}
  \item \emph{Say it:} \emph{kilo}, once more
  \item \emph{Say it:} \emph{do kilo chāvaḷ lāvo.} — Bring two kilos of rice.
  \item \emph{Recall:} say the Marwadi for scales, then the Marwadi for a seer, then say \emph{kilo} again
\end{itemize}

\begin{tcolorbox}[breakable,colback=teal!4,colframe=teal!35!black,title={Before you move on}]
What does \mw{किलो} mean? (\textquotedblleft{}A kilo\textquotedblright{}.) Say it once more.
\end{tcolorbox}
\section[\texorpdfstring{\mw{पाव} (pāv)}{पाव (pāv)}]{\mw{पाव} (pāv) — a quarter (of a seer)}

\label{lesson:MW-C200-pav}



\begin{quote}
Before the new one: say the Marwadi for a seer, then the Marwadi for a kilo.
\end{quote}

\subsection*{You'll want to know: \mw{पाव}}
\textbf{\mw{पाव}} — \emph{pāv} — \textquotedblleft{}a quarter (of a seer)\textquotedblright{}.

\begin{quote}
\mw{एक} \mw{पाव} \mw{घी} \mw{दो।}
\end{quote}

\begin{quote}
\emph{ek pāv ghī do.}
\end{quote}

\begin{quote}
Give me a quarter of ghee.
\end{quote}

A masculine noun.

\begin{cousinweb}
\textbf{\mw{पाव}} (\emph{pāv}) is inherited from Sanskrit \emph{pāda}, \textquotedblleft{}foot; a fourth part\textquotedblright{}; CDIAL lists Old Marwari \emph{pāva}, \textquotedblleft{}foot\textquotedblright{}, and Hindi \emph{pāu}, \textquotedblleft{}quarter\textquotedblright{} (CDIAL 8056).
\end{cousinweb}

\begin{grammarlens}[title={What you've built}]
One more word for reading notices and messages.
\end{grammarlens}

\subsection*{Your turn}
\begin{itemize}
\raggedright
  \item \emph{Say it:} \emph{pāv}
  \item \emph{Say it:} \emph{pāv}, once more
  \item \emph{Say it:} \emph{ek pāv ghī do.} — Give me a quarter of ghee.
  \item \emph{Recall:} say the Marwadi for a seer, then the Marwadi for a kilo, then say \emph{pāv} again
\end{itemize}

\begin{tcolorbox}[breakable,colback=teal!4,colframe=teal!35!black,title={Before you move on}]
What does \mw{पाव} mean? (\textquotedblleft{}A quarter (of a seer)\textquotedblright{}.) Say it once more.
\end{tcolorbox}
\section[\texorpdfstring{\mw{लीटर} (līṭar)}{लीटर (līṭar)}]{\mw{लीटर} (līṭar) — a litre}

\label{lesson:MW-C200-litar}



\begin{quote}
Before the new one: say the Marwadi for a kilo, then the Marwadi for a quarter.
\end{quote}

\subsection*{You'll want to know: \mw{लीटर}}
\textbf{\mw{लीटर}} — \emph{līṭar} — \textquotedblleft{}a litre\textquotedblright{}.

\begin{quote}
\mw{एक} \mw{लीटर} \mw{दूध} \mw{दो।}
\end{quote}

\begin{quote}
\emph{ek līṭar dūdh do.}
\end{quote}

\begin{quote}
Give me a litre of milk.
\end{quote}

A masculine noun.

\begin{cousinweb}
\textbf{\mw{लीटर}} (\emph{līṭar}) is borrowed from English \emph{litre}.
\end{cousinweb}

\begin{grammarlens}[title={What you've built}]
One more word for reading notices and messages.
\end{grammarlens}

\subsection*{Your turn}
\begin{itemize}
\raggedright
  \item \emph{Say it:} \emph{līṭar}
  \item \emph{Say it:} \emph{līṭar}, once more
  \item \emph{Say it:} \emph{ek līṭar dūdh do.} — Give me a litre of milk.
  \item \emph{Recall:} say the Marwadi for a kilo, then the Marwadi for a quarter, then say \emph{līṭar} again
\end{itemize}

\begin{tcolorbox}[breakable,colback=teal!4,colframe=teal!35!black,title={Before you move on}]
What does \mw{लीटर} mean? (\textquotedblleft{}A litre\textquotedblright{}.) Say it once more.
\end{tcolorbox}
